\documentclass[10pt,a4paper]{amsart}
\usepackage[margin=19mm]{geometry}
\usepackage{amsmath,amssymb,mathtools}
\usepackage[scr=rsfs,cal=euler]{mathalfa}
\usepackage{xfrac}
\usepackage[unicode,hidelinks,bookmarks=false]{hyperref}
\newcommand{\ie}{i.e.\ }
\newcommand{\cf}{cf.\ }
\newcommand{\resp}{respectively\ }
\newcommand{\Subsection}{\subsection}
\providecommand{\nobreakdash}{\nobreak\hbox{-}}
\setlength{\emergencystretch}{2em}
\multlinegap0pt
\usepackage{amssymb}
\newtheorem{lemma}{Lemma}
\title{Sign changes of the Liouville function in arithmetic progressions}
\author{Kevin Ford, Maksym Radziwi{\l}{\l}}
\date{Source: arXiv:2605.03349v1 (CC BY 4.0)}
\begin{document}
\maketitle

\noindent\emph{Source and licence.} Sign changes of the Liouville function in arithmetic progressions. Version arXiv:2605.03349v1 (CC BY 4.0), distributed by arXiv under the Creative Commons Attribution 4.0 International licence, http://creativecommons.org/licenses/by/4.0/, as recorded in the arXiv licence metadata for this exact version and shown on the abstract page https://arxiv.org/abs/2605.03349v1. Full licence notice: \texttt{licenses/arxiv-2605.03349-CC-BY-4.0.txt}.\par\smallskip

\noindent\textit{Editorial scope.} Counted unit: the article's lemma lem:periodicity (source0.tex lines 134--148). The article's complete original proof of it is delivered with it (source0.tex lines 150--252, one \texttt{proof} environment of 103 lines). No mathematics is written, repaired or re-derived: the statement and the proof are the article's own lines, in the article's own notation, and no retained line is substituted or re-flowed.  Retained uncounted as the source-local context the proof needs: lines 117--122.  Nothing was omitted from any retained span, so the package carries no omission record.  No retained line contains a citation, so the wrapper carries no bibliography and no bibliography entry.  No retained line contains a figure environment, an image-inclusion command or drawing code, so no image is delivered and the package declares no asset dependency.  Editorial formatting only, in addition: an AMS \texttt{amsart} preamble that restates the article's own declarations and macros used by the retained lines, namely the environments \texttt{lemma} on separate counters, together with the standard packages those lines need.  The retained environments print in this document as Lemma 1 in order of appearance, whereas the article prints the same items with the numbering of its own full document.  The complete native source of the article as distributed in the arXiv e-print for this exact version is bundled unchanged as \texttt{sources/199784/source0.tex}.
\begin{lemma}\label{lem:comb}
Let $A$ and $B$ be two subsets of $(\mathbb{Z}/q\mathbb{Z})^{\times}$
with $|A|+|B|>\varphi(q)$.
Then for every $(n,q)=1$ we can find $a \in A$ and $b \in B$ such that
$n \equiv ab \pmod{q}$.
\end{lemma} 

\begin{lemma}\label{lem:periodicity}
Let $\varepsilon > 0$ and let $q$ be 
sufficiently large  with respect to $1/\varepsilon$.
Let $a$ be an integer coprime to $q$,
let $\kappa \in \{-1, +1\}$, and let $N \geq q^{2 + \varepsilon}$.
If for every $n \equiv a \pmod{q}$ with $n \leq N$ we have
$\lambda(n) = \kappa$, then
\[
\lambda(n + q) = \lambda(n)
\]
for all $n \leq N/q - q$ with $(n, q) = 1$, and
$$
\# \Big \{ 1 \leq n < q : (n,q)=1, \lambda(n) = -1 \Big \} = \frac{\varphi(q)}{2} = \# \Big \{ 1 \leq n < q : (n,q)=1, \lambda(n) = 1 \Big \}.
$$
\end{lemma}

\begin{proof}
Set $\eta := \varepsilon/3$.
For an interval $I = [M, M + q)$ write
\[
A_I = \{n \in I : (n, q) = 1,\; \lambda(n) = -1\},
\qquad
B_I = \{n \in I : (n, q) = 1,\; \lambda(n) = +1\}.
\]
Notice that $|A_I| + |B_I| = \varphi(q)$.

\medskip
\noindent\textbf{Step 1.}
\textit{For every interval $I = [M, M+q) \subset [1, q^{1+\eta})$
we have $|A_I| = |B_I| = \varphi(q)/2$.}

\smallskip
Suppose $|B_I| > \varphi(q)/2$ for some
$I \subset [1, q^{1 + \eta})$. Since $q$ is sufficiently large
with respect to $1/\varepsilon$ and $\eta = \varepsilon / 3$,
we can choose a prime $p \leq q^{\eta}$ coprime to $q$.
Since multiplication by $p$ permutes
$(\mathbb{Z}/q\mathbb{Z})^{\times}$,
the set $pB_I := \{pb \bmod q : b \in B_I\}$
satisfies $|pB_I| = |B_I| > \varphi(q)/2$.

By Lemma~\ref{lem:comb} applied to $(B_I, B_I)$:
for any $(c, q) = 1$ there exist
$b_1, b_2 \in B_I$ with $b_1 b_2 \equiv c \pmod{q}$,
giving $\lambda(b_1 b_2) = +1$ and
$b_1 b_2 \leq q^{2 + 2\eta}$.

By Lemma~\ref{lem:comb} applied to $(B_I, pB_I)$:
there exist $b, b' \in B_I$ with
$pbb' \equiv c \pmod{q}$, giving
$\lambda(pbb') = (-1)(+1)(+1) = -1$ and
$pbb' \leq q^{2 + 3\eta} = q^{2 + \varepsilon} \leq N$.

Taking $c = a$, we obtain integers of both signs in the
class $a$ up to $N$, contradicting the hypothesis.
The case $|A_I| > \varphi(q)/2$ is identical
(swap the roles of $A$ and $B$).
Since $|A_I| \leq \varphi(q)/2$ and $|B_I| \leq \varphi(q)/2$
while $|A_I| + |B_I| = \varphi(q)$,
we conclude $|A_I| = |B_I| = \varphi(q)/2$.

\medskip
\noindent\textbf{Step 2.}
\textit{For every interval $I = [M, M+q) \subset [1, N/q]$
we have $|A_I| = |B_I| = \varphi(q)/2$.}

\smallskip
Suppose $|B_I| > \varphi(q)/2$ for some
$I \subset [1, N/q]$.
Let $J = [1, q)$.
Since $J \subset [1, q^{1 + \eta})$,
Step~1 gives $|A_J| = |B_J| = \varphi(q)/2$.

By Lemma~\ref{lem:comb} applied to $(B_I, A_J)$
(with $|B_I| > \varphi(q)/2$ and $|A_J| = \varphi(q)/2$):
there exist $b \in B_I$, $d \in A_J$
with $bd \equiv a \pmod{q}$.
Then $\lambda(bd) = (+1)(-1) = -1$ and
$bd \leq (N/q) \cdot q = N$.

By Lemma~\ref{lem:comb} applied to $(B_I, B_J)$:
there exist $b_1 \in B_I$, $b_2 \in B_J$ with
$b_1 b_2 \equiv a \pmod{q}$.
Then $\lambda(b_1 b_2) = (+1)(+1) = +1$ and
$b_1 b_2 \leq N$.

Again both signs appear in class $a$ up to $N$,
contradicting the hypothesis.
The case $|A_I| > \varphi(q)/2$ is symmetric.
Since $|A_I| \leq \varphi(q)/2$ and $|B_I| \leq \varphi(q)/2$
while $|A_I| + |B_I| = \varphi(q)$,
we conclude $|A_I| = |B_I| = \varphi(q)/2$.

\medskip
\noindent\textbf{Step 3.}
\textit{For every $n \leq N/q - q$ with $(n,q) = 1$
we have $\lambda(n+q) = \lambda(n)$.}

\smallskip
Let $n \leq N/q - q$ with $(n, q) = 1$.
The intervals
$I = [n, n + q)$ and $I' = [n + 1, n + q + 1)$
both lie in $[1, N/q]$, so Step~2 gives
$|A_I| = |A_{I'}| = \varphi(q)/2$.

Passing from $I$ to $I'$ removes $n$ and adds $n + q$.
Since $q$ is prime and $(n, q) = 1$, we have
$(n + q, q) = 1$, so both $n$ and $n + q$ are counted
by $A$ or $B$.
The cardinality changes as
\[
|A_{I'}| = |A_I|
  - \mathbf{1}[\lambda(n) = -1]
  + \mathbf{1}[\lambda(n + q) = -1].
\]
Since $|A_{I'}| = |A_I|$, we obtain
$\mathbf{1}[\lambda(n) = -1] = \mathbf{1}[\lambda(n+q) = -1]$,
which gives $\lambda(n) = \lambda(n + q)$.
\end{proof}

\end{document}
